\section{Introducción}
\label{sec:introduccion}

El problema del buffer acotado está resuelto desde que Dijkstra lo planteó
con semáforos~\cite{Dijkstra1965}. Lo que cambia cuando los participantes
son \emph{procesos} con espacios de direcciones independientes, y no hilos
de un mismo programa, es casi todo lo que lo rodea: el estado compartido
tiene que vivir en un objeto con nombre que sobrevive a cada participante,
las primitivas de sincronización pasan a ser datos dentro de ese objeto en
vez de servicios del lenguaje, ninguna dirección virtual vale en otro
proceso, y nadie libera los recursos cuando un proceso termina. Birrell
describe el modelo de hilos del que se parte y menciona esta alternativa
solo para advertir que es más costosa~\cite{Birrell1989}.

Este proyecto implementa ese sistema sobre memoria compartida POSIX. Un
\code{creator} crea una instancia con nombre y capacidad; productores y
consumidores independientes se conectan, se registran, trabajan y se
retiran; un \code{finisher} ejecuta un cierre ordenado que drena lo
pendiente, espera a que todos se retiren y libera los recursos. Un quinto
ejecutable, \code{verify}, cruza los registros de todos los procesos y actúa
como oráculo de corrección. La ejecución ocurre en una sola máquina
GNU/Linux y no cubre la caída arbitraria de un proceso: las garantías se
aplican a participantes que siguen el protocolo y a un cierre normal. Esa
frontera no es una omisión cómoda, y la sección~\ref{sec:pruebas} documenta
qué ocurre exactamente cuando se cruza.

La sección~\ref{sec:arquitectura} describe el segmento compartido y qué dato
es privado de cada proceso; la~\ref{sec:protocolo}, el pseudocódigo del
registro, la inserción, la extracción y el cierre, con sus invariantes y sus
puntos de efecto; la~\ref{sec:pruebas}, las pruebas mínimas y el oráculo;
la~\ref{sec:resultados}, el experimento de 45 corridas;
la~\ref{sec:lecturas}, la relación con Dijkstra y Birrell; y
la~\ref{sec:decisiones}, cinco decisiones donde había una alternativa
razonable que se descartó con evidencia. La sección~\ref{sec:gestion}
documenta la trazabilidad del trabajo.
